\input _template

\setlanguage{SK}

\Title{Základy počítania uhlov}

\MathcompsLink{zaklady-pocitania-uhlov-1}

\Author{Patrik Bak}

\sec Úvod

Geometria je najvizuálnejšia časť matematiky -- vieme si ju kresliť na papier, v~softwari ako \Link[https://www.geogebra.org/]{GeoGebra} a~môžeme sa tešiť, keď na vlastné oči vidíme, že veci v~nej fungujú. Prečo tomu tak je, však vie byť záhadou. Aby sme tieto tajomstvá pochopili, potrebujeme kúsok po kúsku budovať základňu znalostí a~vizuálnu intuíciu. 

Jedna z~najzákladnejších techník je rozumieť svetu uhlov a~vedieť ich vlastnosti s~úspechom používať v~rôznych typoch úloh. Len málo zaujímavých príkladov uhly nepoužíva. Teória okolo uhlov pritom nie je zložitá, je však dôležité budovať ju pomaly a detailne. Začneme s~najjednoduchšími vlastnosťami, ktoré ani nevyžadujú prácu s kružnicami. Už tieto sa však dajú s~úspechom využiť v~mnohých úlohách.

V~tomto materiáli sa budeme venovať uhlom, ale zámerne vynecháme uhly v~kružniciach, ktoré si zaslúžia separátnu pozornosť. Na chvíľu zabudnime, že kružnice existujú.

Na záver úvodu jedna menšia poznámka o konvencii. V~tomto materiáli uvidíme veľa obrázkov trojuholníkov, v~ktorých je bod~$A$ nakreslený hore. Ukazuje sa, že tento spôsob kreslenia je medzinárodný štandard, najmä na poli matematickej olympiády a podobných súťaží sa to tak zvykne robiť. V~Česku a na~Slovensku je však v~školách stále zvykom kresliť~$C$ hore. V~niektorých krajinách možno vidieť aj~$B$ hore, napr. na Ukrajine. Z~hľadiska riešenia úloh je vždy dobré nakresliť si hore ten bod, pri ktorom obrázok vyzerá čo najsymetrickejšie, ľahšie v~ňom vidieť ďalšie symetrické veci. V~tomto materiáli budú úlohy zadávané tak, že tá symetria bude vidieť viac pri~$A$ hore \SlightSmile.

\sec Základy sveta uhlov

V~tejto sekcii si odvodíme tie najjednoduchšie vlastnosti súvisiace s uhlami. Pevne verím, že kľúčom k~zvládnutiu geometrie je porozumenie vecí od základov a do hĺbky, preto sa týmto jednoduchým vlastnostiam budeme venovať viac.

Základné vlastnosti uhlov, ktoré bežne používame skoro bez rozmýšľania pri riešení úloh:

\Highlight{
    \begitems \style i
    \i Vrcholové uhly sú zhodné: 
       \Image{angles-vertical.pdf}
    \i Súhlasné uhly sú zhodné:
       \Image{angles-corresponding.pdf}
    \i Striedavé uhly sú zhodné:
       \Image{angles-alternate.pdf}
    \i Vedľajšie uhly majú súčet $180^\circ$:
       \Image{angles-supplementary.pdf}
   \enditems
}

Tieto vlastnosti spolu veľmi prirodzene súvisia, menšie cvičenie na rozmyslenie:

\EnvId{E7dY7LX9YobFdS39tfDDR-vrcholove-suhlasne-striedave-uhly}
\Exercise{}{
    Uvedomte si, že:
    \begitems \style a
    \i z~vlastnosti o~vedľajších uhloch vyplýva vlastnosť vrcholových uhlov (a~naopak)
    \i z~ktorejkoľvek dvojice vlastností o~vrcholových, súhlasných, striedavých uhloch vyplýva tá tretia vlastnosť
    \enditems
}{
    \textbf{(a)} Vedľajší uhol $\alpha'$ k~danému uhlu $\alpha$ má veľkosť $180^\circ - \alpha$. Uhol vedľajší k~$\alpha'$ má zasa veľkosť $180^\circ - (180^\circ - \alpha) = \alpha$. Ten je však zároveň vrcholový uhol k~$\alpha$, teda sú vrcholové uhly zhodné.

    \Image{angles-supplementary-derive-vertical.pdf}

    Naopak, predpokladajme vlastnosť o~vrcholových uhloch a~označme po sebe idúce uhly v~priesečníku dvoch priamok $\alpha$, $\beta$, $\gamma$, $\delta$. Zo zhodnosti vrcholových uhlov $\alpha = \gamma$ a~$\beta = \delta$. Súčet všetkých štyroch je plný uhol, teda $\alpha + \beta + \gamma + \delta = 360^\circ$. Po dosadení dostávame $2\alpha + 2\beta = 360^\circ$, čiže $\alpha + \beta = 180^\circ$ -- to je práve vlastnosť vedľajších uhlov.

    \Image{angles-vertical-derive-supplementary.pdf}

    \textbf{(b)} Všetky tri vlastnosti hovoria, že istá dvojica uhlov má rovnakú veľkosť.

    Pozrime sa na trojicu uhlov $\alpha$, $\beta$, $\gamma$ z~obrázka a~na to, čo o~ich vzťahoch tvrdia jednotlivé vlastnosti. Pre dvojicu $\alpha$, $\beta$ pri hornom priesečníku tvrdí vlastnosť vrcholových uhlov $\alpha = \beta$. Pre dvojicu $\alpha$, $\gamma$ tvrdí vlastnosť súhlasných uhlov $\alpha = \gamma$. Konečne pre dvojicu $\beta$, $\gamma$ tvrdí vlastnosť striedavých uhlov $\beta = \gamma$.

    \Image{angles-corresponding-alternate-supplementary-connection.pdf}

    Tri vlastnosti teda tvrdia tri rovnosti medzi tými istými troma uhlami, len cez rôzne dvojice. Akonáhle platia ľubovoľné dve, tretia plynie tranzitivitou: napr. ak vieme (V) a~(S), tak $\alpha = \beta$ a~$\alpha = \gamma$, takže $\beta = \gamma$, čiže (St). Ostatné dvojice analogicky.
}

Nasledujúce tvrdenie všetci poznáme, viete ho však dokázať?

\EnvId{G2fX3226MoyMyLEg7bCt3-sucet-uhlov-v-trojuholniku}
\Theorem{}{
    Súčet uhlov v~trojuholníku je $180^\circ$.
}{
    Veďme cez vrchol~$A$ priamku rovnobežnú s~$BC$. Pretože $AB$ je priečka rovnobežiek, uhol $\beta$ pri~$A$ a~uhol $\beta$ pri~$B$ sú striedavé, teda rovnaké. Analogicky cez priečku $AC$ sú uhly $\gamma$ pri~$A$ a~pri~$C$ striedavé. Uhly $\beta$, $\alpha$, $\gamma$ ležia vedľa seba pozdĺž priamky cez~$A$, teda  $\alpha + \beta + \gamma = 180^\circ.$

    \Image{angles-triangle.pdf}
}

Toto tvrdenie sa rozhodne dá zovšeobecniť. Vieme napríklad, že súčet uhlov v~štvoruholníku je $360^\circ$. Ako je to v 5-uholníku? A ako v~67-uholníku? Odpoveďou je nasledujúce tvrdenie:

\EnvId{H7yQvLt3VxyGlELBTgAJh-sucet-uhlov-konvexneho-mnohouholnika}
\Theorem{}{
    Nech $n \ge 3$ je prirodzené číslo. Súčet uhlov v~konvexnom $n$-uholníku je rovný $(n-2) \cdot 180^\circ$.
}{
    \NamedProof{Dôkaz 1 (matematická indukcia).} Základ $n=3$: súčet uhlov trojuholníka je $180^\circ = (3-2)\cdot 180^\circ$. Indukčný krok: nech $A_1 A_2 \ldots A_{n+1}$ je konvexný $(n+1)$-uholník. Uhlopriečkou $A_1 A_n$ ho rozdelíme na trojuholník $A_1 A_n A_{n+1}$ a~konvexný $n$-uholník $A_1 A_2 \ldots A_n$: uhly pri vrcholoch $A_2, \ldots, A_{n-1}$ patria celé $n$-uholníku, uhol pri $A_{n+1}$ patrí celý trojuholníku a~uhly pri $A_1$, $A_n$ sa rozdelia medzi oba útvary tak, aby ich časti dali pôvodné vnútorné uhly $(n+1)$-uholníka. Súčet uhlov $(n+1)$-uholníka je teda
    $$
    \underbrace{180^\circ}_{\text{trojuholník}} + \underbrace{(n-2) \cdot 180^\circ}_{n\text{-uholník}} = (n-1) \cdot 180^\circ = \bigl((n+1) - 2\bigr) \cdot 180^\circ.
    $$

    \Image{angles-polygon.pdf}

    \NamedProof{Dôkaz 2 (vejárová triangulácia).} Veďme z~vrcholu $A_1$ uhlopriečky do každého nepriliehajúceho vrcholu $A_3, A_4, \ldots, A_{n-1}$. Vznikne $n-2$ trojuholníkov pokrývajúcich vnútro $n$-uholníka bez prekrytia. Pri každom vrchole $A_k$ sa uhly susedných trojuholníkov pri $A_k$ presne poskladajú na vnútorný uhol $n$-uholníka pri $A_k$, takže súčet uhlov všetkých trojuholníkov sa rovná súčtu vnútorných uhlov $n$-uholníka, čiže $(n-2)\cdot 180^\circ$.

    \Image{angles-polygon-fan.pdf}

    \Remark Dôkaz~2 je \uv{rozbalený} dôkaz~1: indukcia odrezáva trojuholníky jeden po druhom, tu ich vidíme všetky naraz.

    \NamedProof{Dôkaz 3 (vnútorný bod).} Zvoľme bod $O$ vo vnútri $n$-uholníka a~spojme ho s~každým vrcholom. Vznikne $n$ trojuholníkov s~celkovým súčtom uhlov $n\cdot 180^\circ$. Uhly trojuholníkov pri každom vrchole $A_k$ spolu tvoria vnútorný uhol $n$-uholníka pri~$A_k$; uhly trojuholníkov pri~$O$ spolu tvoria plný uhol $360^\circ$. Teda súčet vnútorných uhlov $n$-uholníka je
    $$
    n\cdot 180^\circ - 360^\circ = (n-2)\cdot 180^\circ.
    $$

    \Image{angles-polygon-interior.pdf}

    \Remark{nekonvexné mnohouholníky} Vzorec $(n-2)\cdot 180^\circ$ platí aj pre nekonvexné mnohouholníky, ak všetky vnútorné uhly meriame smerom dovnútra -- teda aj uhly väčšie než $180^\circ$ rátame ako také.

    Vejárová triangulácia z~jedného vrcholu fungovať nemusí -- niektoré uhlopriečky môžu prechádzať mimo mnohouholníka. Platí však, že každý jednoduchý $n$-uholník sa dá rozdeliť na $n-2$ trojuholníkov (nie nutne vejárovo), čo dôkaz zachráni. Dôvod je indukcia cez trojuholník tvorený nejakým vrcholom $V$ a~jeho susedmi $U$, $W$, ktorý leží celý vo vnútri nášho $n$-uholníka. Odrezaním $V$ uhlopriečkou $UW$ dostaneme $(n-1)$-uholník; opakovaním teda $n-2$ trojuholníkov. Dá sa dokázať, že takýto trojuholník vždy existuje (dokonca dva), viď \Link[https://en.wikipedia.org/wiki/Two_ears_theorem]{veta o~dvoch ušiach}.

    K~dôkazu~3: funguje len ak existuje vnútorný bod $O$ viditeľný zo všetkých vrcholov (t. j. úsečky $OA_i$ ležia celé vo vnútri). Pre niektoré nekonvexné mnohouholníky taký bod $O$ nemusí existovať (skúste taký nakresliť) -- v~takom prípade sa dôkaz~3 nedá priamo zachrániť a~treba siahnuť po triangulácii popísanej v~predošlom odstavci.
}

Teraz uveďme ešte dve jednoduché, ale užitočné pomocné tvrdenia:

\EnvId{vUBGcz0IXvF69ZPtrH_4F-sucet-uhlov-na-obrazku}
\Exercise{}{
    Dokážte, že súčet uhlov na obrázku je rovný $180^\circ$.

    \Image{angles-complementary-parallel-statement.pdf}
}{
    Označme $\alpha$ uhol zvieraný priečkou $AB$ s~rovnobežkami na strane $B$ (ako na obrázku). Súhlasné uhly $\alpha$ pri~$A$ a~$\alpha$ pri~$B$ sú zhodné. Pri vrchole $A$ sú $\alpha$ a~$\beta$ vedľajšie, teda
    $$
    \beta + \alpha = 180^\circ.
    $$

    \Image{angles-complementary-parallel-solution.pdf}
}

\EnvId{EMfGZzO5Ro3VypJ7caLES-uhol-s-otaznikom-na-obrazku}
\Exercise{}{
    Dokážte, že uhol s~otázníkom na obrázku je rovný $\alpha+\beta$.

    \Image{angles-complementary-in-triangle.pdf}
}{
    Zo súčtu uhlov v~trojuholníku $|\angle ACB| = 180^\circ - \alpha - \beta$. Vonkajší uhol pri $C$ je vedľajší k~$|\angle ACB|$:
    $$
    180^\circ - |\angle ACB| = \alpha + \beta.
    $$

    \Image{angles-complementary-in-triangle.pdf}
}

Posledné dve cvičenia možno vyzerali veľmi triviálne. Realita však je, že v~praktických úlohách sú veľmi užitočné -- keď máme obrovský netriviálny obrázok a~máme previesť za sebou veľký počet uhlových operácií, tak je naozaj veľmi výhodné si čo i len kúsok tej práce zjednodušiť -- v~prípade rovnobežiek nemusíme uvažovať pomocný súhlasný uhol a~v~prípade trojuholníka nemusíme pracovať s~$180^\circ$.

\sec Základy sveta dĺžok

Zatiaľ sme fungovali iba vo svete uhlov a~vôbec sme neriešili dĺžky. Geometria však začne byť zaujímavá, keď tieto svety začneme prepájať. Základom pre nás bude {\em zhodnosť trojuholníkov}. Pripomeňme si kritériá.

\Highlight{
    Vety o~zhodnosti trojuholníkov:
    \begitems \style i
    \i Veta $sss$: dva trojuholníky sú zhodné, ak sa zhodujú vo všetkých troch stranách.
    \i Veta $sus$: dva trojuholníky sú zhodné, ak sa zhodujú v~dvoch stranách a~uhle, ktorý tieto strany zvierajú.
    \i Veta $usu$: dva trojuholníky sú zhodné, ak sa zhodujú v~jednej strane a~nejakých dvoch uhloch.
    \i Veta $Ssu$: dva trojuholníky sú zhodné, ak sa zhodujú v~dvoch stranách a~uhle oproti dlhšej z~týchto strán.
    \enditems
}

Zdôraznime, že pri vete $sus$ je naozaj dôležité, že zhodný uhol je ten zvieraný dvoma stranami -- bez tohto predpokladu zhodnosť nemusí fungovať, viď obrázok. Na ňom máme dva trojuholníky $ABC$ a $A'B'C'$, pre ktoré platí $|BC|=|B'C'|$, $|AC|=|A'C'|$, a $\beta=\beta'$, avšak zjavne nie sú zhodné.

\Image{angles-sas-warning.pdf}

Veta $Ssu$ by bola naša záchrana -- tu ju však aplikovať nevieme, lebo strana, oproti ktorej je uhol $\beta$, teda $|AC|$, nie je najdlhšia, keďže je zjavne kratšia než $|BC|$ \SmileWithTear.

Než pôjdeme ďalej, uvedomme si, čo tieto vety znamenajú. Za mňa dobrý pohľad na zhodnosť je takýto: chceme zostrojiť trojuholník, keď máme dané nejaké tri jeho elementy -- budú všetky zostrojiteľné trojuholníky zhodné?

\EnvId{YVYLXXXSL6OFFUuABkNDD-veta-sss-o-zhodnosti}
\Example{}{
    Presvedčte sa o~platnosti vety $sss$ o~zhodnosti. 
}{
    Predstavme si, že máme dané tri úsečky dĺžok $a$, $b$, $c$, z~ktorých vieme zostrojiť trojuholník. Bez ujmy na všeobecnosti predpokladajme, že $a=\max\{a,b,c\}$. Začnime konštruovať trojuholník úsečkou~$BC$ dĺžky~$a$. Následne  zostrojíme dve kružnice: (a) kružnicu so stredom v~$B$ a~polomerom $c$; (b) kružnicu so stredom v~$C$ a~polomerom $b$. Vďaka $a \ge b$ resp. $a \ge c$ obe tieto kružnice pretínajú úsečku~$BC$. V~prípade $a=b+c$ sa pretínajú práve v~jednom bode na nej (čo nezodpovedá trojuholníku), v~prípade $a>b+c$ sa nepretínajú vôbec, a~v~prípade $a<b+c$ sa pretínajú v~dvoch bodoch $A$ a~$A'$.

    \Image{angles-sss-proof.pdf}

    Na konci si uvedomme, že trojuholníky $ABC$ a~$A'BC$ sú však zrejme vzájomné zrkadlové obrazy podľa $BC$. Taktiež ak by sme začali inou úsečkou, tak vytvoríme iba posunutie/otočenie tejto konfigurácie.   

    \Remark Uvedomme si, že sme vlastne po ceste dokázali~trojuholníkovú nerovnosť. Tá sa zvyčajne formuluje ako: súčet ktorýchkoľvek dvoch strán trojuholníka je väčší než tretia. My sme dokázali, že súčet dvoch najkratších je väčší než najdlhšia. To zrejme stačí, keďže súčet najdlhšej a~ktorejkoľvek inej je väčší než tá posledná. 
}

Podobne vieme zdôvodniť ďalšie kritériá: 

\EnvId{vMyC0sIAGJZq-NeOnkTWW-veta-sus-o-zhodnosti}
\Exercise{}{
    Presvedčte sa o~platnosti vety $sus$ o~zhodnosti.
}{
    Nech sú dané strany $b = |AC|$, $c = |AB|$ a~uhol $|\angle BAC| = \alpha$ medzi nimi. Zostrojme vrchol~$A$ a~úsečku $AB$ dĺžky~$c$ -- tá je daná až na zhodné zobrazenie. Bod $C$ musí ležať na polpriamke z~$A$ zvierajúcej s~$AB$ uhol~$\alpha$ a~zároveň vo vzdialenosti~$b$ od~$A$. Polpriamky zvierajúce s~$AB$ uhol $\alpha$ sú dve (po jednej na každej strane $AB$): na každej leží práve jeden bod vo vzdialenosti $b$ od~$A$, teda dostaneme body $C$ a~$C'$. Trojuholníky $ABC$ a~$ABC'$ sú zrkadlové obrazy podľa $AB$, teda zhodné.

    \Image{angles-sas-proof.pdf}
}

\EnvId{PmP9xfoxbBUEnYfG89CCT-veta-usu-o-zhodnosti}
\Exercise{}{
    Presvedčte sa o~platnosti vety $usu$ o~zhodnosti.
}{
    V~prvom rade si uvedomme, že nezáleží na tom, ktoré dva uhly sú zhodné -- ak sa dva trojuholníky zhodujú v~dvoch uhloch, tak tretí je určený jednoznačne, keďže všetky tri majú súčet uhlov $180^\circ$.

    Pre účely nášho dôkazu uvážme teda, že sú zhodné práve uhly priliehajúce zhodnej strane. Nech je to uhol $\beta$ pri vrchole $B$, $\gamma$ pri vrchole $C$ a~strana $a = |BC|$ medzi nimi. Zostrojme úsečku $BC$ dĺžky~$a$ -- tá je daná až na zhodné zobrazenie. Bod $A$ musí ležať na polpriamke z~$B$ zvierajúcej s~$BC$ uhol $\beta$ a~zároveň na polpriamke z~$C$ zvierajúcej s~$CB$ uhol $\gamma$. Súčet vnútorných uhlov pri $B$ a~$C$ v~trojuholníku je menší než $180^\circ$, čiže $\beta + \gamma < 180^\circ$, a~preto nie sú obe polpriamky rovnobežné -- pretnú sa v~jedinom bode~$A$. Trojuholník $ABC$ je teda určený jednoznačne (až na zhodné zobrazenia).

    \Image{angles-asa-proof.pdf}
}

Menej známe tvrdenie $Ssu$ je už máličko ťažšie, preto k~nemu budú aj návody \SlightSmile.

\EnvId{G0TTdrAZrHu-Oip2mtDlP-zhodnost-ssu}
\Problem{0}{zhodnosť $Ssu$}{
    Presvedčte sa o~platnosti vety $Ssu$ o~zhodnosti. Navyše, uvedomte si, kde by sa konštrukcia pokazila, keby uhol nebol voči väčšej zo strán -- čo by sa stalo, keby boli obe strany rovnako dlhé?
}{
    Začnite konštrukciu konštrukciou kratšej úsečky. Ako pokračuje konštrukcia po jej zostrojení?
}{
    Po zostrojení kratšej úsečky zostrojíme polpriamku pod naším daným zhodným uhlom. Zostáva posledný krok konštrukcie. Rozmyslite si, kedy dostaneme žiaden, kedy jeden a~kedy dva vyhovujúce trojuholníky.
}{
    Nech sú dané strany $a = |BC|$, $b = |AC|$ a~uhol $|\angle ABC| = \beta$ oproti dlhšej strane $b$. Začneme zostrojením kratšej úsečky $BC$ dĺžky~$a$. Bod $A$ musí ležať na polpriamke z~$B$ zvierajúcej s~$BC$ uhol $\beta$ a~zároveň vo vzdialenosti $b$ od~$C$, čiže na kružnici $k$ so stredom~$C$ a~polomerom $b$.

    Vzdialenosť bodu~$B$ od stredu kružnice je $|BC| = a$. Pretože $a < b$, bod~$B$ leží \textit{vnútri} kružnice $k$. Polpriamka z~bodu~$B$ tak začína vnútri kružnice a~pokračuje smerom von, takže kružnicu pretne práve raz. Bod $A$ je teda určený jednoznačne, a~trojuholník $ABC$ je takisto jednoznačný (až na zhodné zobrazenie).

    \Image{angles-Ssa-proof-one-solution.pdf}

    Pozrime sa, prečo je podmienka $b > a$ podstatná. Ak by bolo $a = b$, bod~$B$ by ležal priamo \textit{na} kružnici~$k$. Polpriamka z~$B$ by ju pretínala v~bode~$B$ samotnom (čo nedáva trojuholník) a~v~práve jednom ďalšom bode -- ten je hľadaným~$A$. Konštrukcia teda stále vedie k~jednoznačnému (rovnoramennému) trojuholníku, takže veta $Ssu$ formálne platí aj v~tomto hraničnom prípade. Pod $Ssu$ ho však samostatne nezaraďujeme: ak vieme, že trojuholník je rovnoramenný so $|BC|=|AC|$ a~poznáme jeden uhol $\beta$, tak ostatné dva uhly sú už vďaka rovnoramennosti určené (oba pri základni majú veľkosť~$\beta$). Tá istá konfigurácia je preto pokrytá vetami $sus$ či $sss$ a~$Ssu$ k~nej nepridáva žiadnu novú informáciu.

    \Image{angles-Ssa-proof-isosceles.pdf}

    Ak by bolo $b < a$, bod $B$ by ležal \textit{zvonku} kružnice $k$. Polpriamka by ju mohla pretnúť v~dvoch bodoch -- vznikli by tak dva rôzne (nezhodné) trojuholníky.

    \Image{angles-Ssa-proof-two-solutions.pdf}

    \Remark Dôkaz je možné spraviť aj tak, že najprv zostrojíme dlhšiu úsečku. To však vyžaduje znalosť množiny bodov nad pevnou úsečkou majúcich pevný uhol -- k~tomuto sa dostaneme v~inom materiáli.
}

Na našej geometrickej ceste sa priebežne stretneme so všetkými týmito tvrdeniami. Nateraz si však ukážme nejaké konkrétne aplikácie. Začneme tým najzrejmejším tvrdením, ktoré sa ale tiež sluší a patrí dokázať:

\EnvId{GMR2XuASqafLT4gCEpuRI-rovnoramenny-trojuholnik-zhodne-uhly}
\Theorem{rovnoramenný trojuholník}{
    Dokážte, že ak pre trojuholník~$ABC$ platí $|AB|=|AC|$, tak $|\angle ABC|=|\angle ACB|$. Dokážte tiež obrátenú implikáciu (teda že z~rovnosti uhlov vyplýva rovnosť dĺžok).
}{
    \textit{Priama implikácia.} Predpokladajme $|AB| = |AC|$. Nech $M$ je stred $BC$. Trojuholníky $ABM$ a~$ACM$ sú zhodné podľa $sss$ ($|AB|=|AC|$ z~predpokladu, $|BM|=|CM|$ zo~stredu, spoločná strana $AM$), odkiaľ $|\angle ABC| = |\angle ACB|$. Tento dôkaz funguje aj s~inými voľbami bodu $M$. Ak za $M$ vezmeme pätu kolmice z~$A$ na $BC$, dostaneme spoločnú stranu $AM$, $|AB|=|AC|$ a~pravý uhol pri~$M$, ktorý leží oproti najdlhšej strane (prepone) $AB$ resp.~$AC$ -- aplikujeme vetu $Ssu$. Ak za $M$ vezmeme pätu osi vnútorného uhla pri~$A$, máme $|AB|=|AC|$, zhodný uhol $|\angle BAM|=|\angle CAM|$ a~spoločnú stranu $AM$ medzi nimi -- aplikujeme vetu $sus$.

    \Image{angles-isosceles-triangle.pdf}

    \textit{Obrátená implikácia.} Predpokladajme $|\angle ABC| = |\angle ACB| = \beta$. Nech $M$ je päta kolmice z~$A$ na $BC$. Podľa vety $usu$  je $\triangle ABM \cong \triangle ACM$, keďže máme dva zhodné uhly $|\angle ABM| = |\angle ACM|$ a $|\angle AMB| = |\angle AMC| = 90^\circ$ a~spoločnú stranu~$AM$, teda $|AB| = |AC|$. Podobne by sme za $M$ mohli zvoliť pätu osi vnútorného uhla pri~$A$ -- zo zhodných uhlov $|\angle BAM|=|\angle CAM|$ a~$|\angle ABM|=|\angle ACM|$ a~zo spoločnej strany $AM$ znova aplikujeme vetu $usu$. Voľba $M$ ako stredu $BC$ však tu nepomáha: dostaneme síce $|BM|=|CM|$, $AM$ spoločnú a~$|\angle ABM|=|\angle ACM|=\beta$, no uhol $\beta$ leží oproti $AM$, ktorá môže byť kratšia než $BM$ (pri tupom uhle pri~$A$), takže ani vetu $Ssu$ vo všeobecnosti nevieme použiť -- poučenie je, že aj keď máme správny bod, tak na jeho presnej definícii záleží (to ešte veľakrát uvidíme).

    \Remark Iný vtipný dôkaz je založený na tom, že dokážeme $\triangle ABC \cong \triangle ACB$ (všimnite si rôzne poradie vrcholov). V~prípade, že poznáme rovnaké strany, tak použijeme vetu $sus$ alebo dokonca $sss$ -- tá nám následne dá zhodné uhly. V~prípade, že poznáme rovnaké uhly, to zas bude veta $usu$, a~zhodnosť nám dá rovnaké strany.
}

Vyskúšajte si exaktne dokázať tieto jednoduché poznatky, to už bude jednoduchšie: 

\EnvId{rKssiyhl1D8AnznvYHFaN-uhly-rovnostranneho-trojuholnika}
\Exercise{}{
    Rovnostranný trojuholník má tri zhodné uhly rovné $60^\circ$.
}{
    V~$\triangle ABC$ s~$|AB| = |AC| = |BC|$ dáva veta o~rovnoramennom trojuholníku $|\angle ABC| = |\angle ACB|$ (z~$|AB|=|AC|$); analogicky z~$|AB|=|BC|$ je $|\angle BAC| = |\angle BCA|$. Všetky tri uhly sú teda zhodné, a~zo súčtu $180^\circ$ je každý $60^\circ$.

    \Image{angles-equilateral.pdf}
}

\EnvId{Dh7EmG5yYRQD44zhd5YD6-rovnoramenny-pravouhly-trojuholnik}
\Exercise{}{
    Rovnoramenný pravouhlý trojuholník má tri uhly rovné $90^\circ$, $45^\circ$, $45^\circ$.
}{
    Nech je pravý uhol pri vrchole~$A$, teda $|\angle BAC| = 90^\circ$, a~odvesny sú zhodné, $|AB| = |AC|$. Podľa vety o~rovnoramennom trojuholníku platí $|\angle ABC| = |\angle ACB|$. Označme túto spoločnú veľkosť $\beta$. Zo súčtu uhlov v~trojuholníku
    $$
    90^\circ + \beta + \beta = 180^\circ,
    $$
    odkiaľ $2\beta = 90^\circ$, čiže $\beta = 45^\circ$.

    \Image{angles-isosceles-right-triangle.pdf}
}

Veľmi šikovná vec použiteľná v~neľahkých úlohách je nasledovná úloha. Jeden možný dôkaz je cez trigonometriu. My to ale chceme pekne geometricky.

\EnvId{PVssK3ww1Ww4_Lm3f4PrT-pravouhly-trojuholnik-dvojnasobna-prepona}
\Problem{0}{}{
    Pravouhlý trojuholník má zvyšné dva uhly rovné $60^\circ$ a~$30^\circ$ práve vtedy, keď je jeho prepona dvakrát dlhšia než kratšia odvesna.
}{
    Povedzme, že $AB$ je prepona a~je dvakrát dlhšia než odvesna $BC$. Trikom je uvážiť bod $B'$ taký, že $C$ je stred úsečky $BB'$.
}{
    Nech náš trojuholník má pravý uhol pri vrchole~$C$.

    \textit{($\Rightarrow$)} Predpokladajme $|AB| = 2|BC|$. Nech $B'$ je taký bod, že $C$ je stredom úsečky $BB'$. Potom $|B'C| = |BC|$ a~$C$ leží medzi $B$, $B'$, takže $|\angle B'CA|$ je vedľajší k~$|\angle BCA| = 90^\circ$, čiže tiež $90^\circ$. Platí $\triangle CBA \cong \triangle CB'A$ z~vety $sus$, keďže uhly pri~$C$ sú oba pravé, strana~$AC$ je spoločná, a $|BC|=|B'C|$.

    \Image{angles-90-60-30.pdf}

    Pretože $C$ je stredom $BB'$, je $|BB'| = 2|BC|$, a~zo zadania $|AB| = 2|BC|$. Spolu
    $$
    |BB'| = |AB| = |AB'|,
    $$
    takže $\triangle ABB'$ je rovnostranný a~všetky jeho uhly sú $60^\circ$. Špeciálne $|\angle ABC| = |\angle ABB'| = 60^\circ$. Zo súčtu uhlov v~$\triangle ABC$ dopočítame $|\angle BAC| = 180^\circ - 90^\circ - 60^\circ = 30^\circ$.

    \textit{($\Leftarrow$)} Nie je ťažké rozmyslieť si, že úvahy z~predošlého odstavca vieme ľahko obrátiť -- kľúčovú zhodnosť $\triangle CBA \cong \triangle CB'A$ tentokrát dostaneme z~vety $usu$.

    \Remark Iné riešenie je uvážiť, že podľa Tálesovej vety je stred~$O$ kružnice opísanej trojuholníku $ABC$ zároveň stredom prepony $AB$. K~tomuto riešeniu sa podrobnejšie vrátime, keď sa budeme baviť o~kružniciach.
}

\sec Osi

Zo školy iste poznáme dve osi: os úsečky a os uhla. V~tejto sekcii si uzrejmíme známe veci a~dôležité vlastnosti, ktoré budeme priebežne používať. 

\secc Os strany 

Os strany vieme definovať ako priamku, ktorá je kolmá na danú úsečku a~prechádza jej stredom. Iná definícia je, že je to množina bodov, ktoré majú rovnakú vzdialenosť od krajných bodov našej úsečky. Sú tieto definície ale ekvivalentné? Nuž, to nie je evidentné, zdôvodníme si to.

\EnvId{kTzIqbV0OaJSgPvurjL5y-os-usecky}
\Theorem{os úsečky}{
    Množina bodov, ktoré majú rovnakú vzdialenosť od krajných bodov úsečky $AB$, je kolmica prechádzajúca jej stredom~$M$.
}{
    \textit{($\Rightarrow$)} Nech bod~$X$ leží na kolmici na $AB$ vedenej stredom~$M$. Trojuholníky $AMX$ a~$BMX$ sú zhodné podľa vety $sus$: $|MA|=|MB|$ (stred), spoločná strana $MX$ a~zhodné pravé uhly pri~$M$. Odtiaľ $|XA|=|XB|$.

    \Image{angles-perpendicular-bisector-forward.pdf}

    \textit{($\Leftarrow$)} Predpokladajme, že bod~$X$ spĺňa $|XA|=|XB|$, ale neleží na kolmici na $AB$ cez~$M$. Ak $X$ leží na priamke $AB$, je zrejmé, že jediný bod tejto priamky rovnako vzdialený od~$A$ aj~$B$ je práve stred~$M$.

    Predpokladajme teda, že $X$ na priamke $AB$ neleží. Bez ujmy na všeobecnosti nech je naša kolmica zvislá a~$X$ leží naľavo od nej. Označme $Y$ priesečník úsečky $XB$ s~touto kolmicou. Podľa už dokázanej implikácie platí $|YA|=|YB|$.

    \Image{angles-perpendicular-bisector-converse.pdf}

    V~rovnoramennom $\triangle XAB$ ($|XA|=|XB|$) je $|\angle XAB|=|\angle XBA|$. Keďže $Y$ leží na úsečke $XB$, polpriamky $BX$ a~$BY$ splývajú, takže $|\angle XBA|=|\angle YBA|$. Avšak v~rovnoramennom $\triangle YAB$ ($|YA|=|YB|$) je $|\angle YBA|=|\angle YAB|$. Spolu
    $$
    |\angle XAB|=|\angle XBA|=|\angle YBA|=|\angle YAB|.
    $$
    To je však nezmysel: bod~$Y$ leží vnútri úsečky $XB$, preto polpriamka $AY$ smeruje do vnútra trojuholníka $XAB$ a~leží striktne vnútri uhla $\angle XAB$, čiže $|\angle YAB|<|\angle XAB|$.
}

Najzákladnejšia veta zahŕňajúca osi strán je samozrejme nasledovné tvrdenie:

\EnvId{FNKSPjq6yM0S5N_PL53GL-existencia-stredu-opisanej-kruznice}
\Theorem{existencia~$O$}{
    Osi strán ľubovoľného trojuholníka sa pretínajú v~jednom bode. Tento je stredom kružnice opísanej danému trojuholníku.
}{
    Nech $ABC$ je trojuholník. Os úsečky $AB$ je kolmá na $AB$ a~os úsečky $AC$ je kolmá na $AC$; tieto dve kolmice sú rovnobežné iba v~prípade $AB \parallel AC$, čo zrejme neplatí. Preto sa osi úsečiek $AB$ a~$AC$ pretínajú v~jednom bode -- označme ho~$O$.

    Podľa predošlej vety platí $|OA|=|OB|$ (lebo $O$ leží na osi $AB$) a~$|OA|=|OC|$ (lebo $O$ leží na osi $AC$). Dokopy to dáva $|OB|=|OC|$, čo nám zasa spätne dáva, že $O$ leží aj na osi úsečky $BC$. Všetky tri osi teda prechádzajú bodom~$O$.

    \Image{angles-circumcenter.pdf}

    \Remark{1} Bod $O$ spĺňa $|OA|=|OB|=|OC|$, takže je stredom kružnice prechádzajúcej cez všetky tri vrcholy -- teda kružnice opísanej trojuholníku $ABC$.

    \Remark{2} Označenie $O$ pre stred opísanej kružnice je medzinárodný štandard. V~Česku a na Slovensku sa zvykne používať aj $S$ (zo slova \uv{stred}).
}

\secc Os uhla

Podobne ako pri osi strany sa najprv zamyslíme nad definíciou. Majme uhol $XAY$. Pod jeho \textit{osou} (presnejšie \textit{vnútornou osou}) rozumieme polpriamku~$AZ$ takú, že $|\angle XAZ|=|\angle ZAY|$. Bod~$Z$ je teda zvolený tak, že polpriamka~$AZ$ rozdelí uhol $XAY$ na dva zhodné uhly. Voľne sa za os uhla pokladá aj~celá priamka určená touto polpriamkou.

\Image{angles-bisector-definition.pdf}

Rovnako ako pri osi strany sa táto definícia často zamieňa s~tvrdením, že os uhla je množina vnútorných bodov uhla, ktoré majú rovnakú vzdialenosť od jeho ramien. Aj~tu si to zaslúži dôkaz.

\EnvId{PoEX7OH-T0op5IPL9DEhz-os-uhla}
\Theorem{os uhla}{
    Bod~$Z$ vnútri uhla $XAY$ má rovnakú vzdialenosť od priamok $AX$ a~$AY$ práve vtedy, keď $|\angle XAZ|=|\angle ZAY|$.
}{
    \textit{($\Rightarrow$)} Nech polpriamka~$AZ$ delí uhol $XAY$ na dva zhodné uhly. Označme $X'$, $Y'$ päty kolmíc spustených z~$Z$ na priamky $AX$, $AY$. Pravouhlé trojuholníky $AX'Z$ a~$AY'Z$ sú zhodné podľa vety~$usu$: spoločná prepona $AZ$, zhodné uhly pri~$A$ a~zhodné pravé uhly pri~$X'$, $Y'$. Odtiaľ $|ZX'|=|ZY'|$, čo sú práve vzdialenosti bodu~$Z$ od priamok $AX$ a~$AY$.

    \Image{angles-bisector-equidistant.pdf}

    \textit{($\Leftarrow$)} Nech naopak $|ZX'|=|ZY'|$, kde $X'$, $Y'$ sú opäť päty kolmíc z~$Z$ na priamky $AX$, $AY$. Pravouhlé trojuholníky $AX'Z$ a~$AY'Z$ sú zhodné podľa vety~$Ssu$: spoločná prepona $AZ$, zhodné odvesny $|ZX'|=|ZY'|$ a~zhodné pravé uhly pri~$X'$, $Y'$ oproti tejto prepone (najdlhšej strane). Odtiaľ $|\angle X'AZ|=|\angle Y'AZ|$, čo je presne $|\angle XAZ|=|\angle ZAY|$.
}

Než si túto vetu naplno užijeme, zaveďme ešte pojem \textit{vonkajšej osi uhla}. Vonkajšou osou uhla $XAY$ rozumieme zjednotenie osí dvoch uhlov vedľajších k~$XAY$ -- označme si ich $XAY'$ a~$X'AY$, kde $X'$, $Y'$ ležia na opačných polpriamkach k~$AX$, $AY$. Osi týchto dvoch vedľajších uhlov sú opačné polpriamky vychádzajúce z~bodu $A$; ich zjednotenie je teda priamka prechádzajúca~$A$.

\Image{angles-bisector-external-definition.pdf}

Než pôjdeme ďalej, dokážme si túto základnú vlastnosť:

\EnvId{g42rGrZiZDuBQcCBvCKh8-vnutorna-a-vonkajsia-os-uhla}
\Exercise{}{
    Dokážte, že vnútorná os uhla je kolmá na jeho vonkajšiu os.
}{
    Označme $\beta$ veľkosť polovice uhla $XAY$ a~$\alpha$ veľkosť polovice uhla vedľajšieho k~$XAY$, ako na obrázku. Súčet všetkých štyroch uhlov na priamke~$AY$ je
    $$
    \beta+\beta+\alpha+\alpha=180^\circ,
    $$
    z~čoho $\alpha+\beta=90^\circ$. Lenže $\alpha+\beta$ je práve uhol medzi vnútornou a~vonkajšou osou, takže sú kolmé.

    \Image{angles-bisectors-perpendicular.pdf}
}

Vonkajšie osi majú prirodzene tiež vlastnosť, že ich body sú rovnako vzdialené od oboch ramien -- veď sú to stále len osi nejakých uhlov. Ak teda hľadáme množinu všetkých bodov, ktoré majú rovnakú vzdialenosť od dvoch nerovnobežných \textit{priamok} (t.j. nie len polpriamok), dostaneme \textit{zjednotenie} vnútornej a~vonkajšej osi uhla, ktorý tieto priamky zvierajú.

\Image{angles-two-lines-bisectors.pdf}

Teraz prejdime na najznámejší dôsledok súvisiaci s~osami uhlov. Týka sa trojuholníka a~musíme v~ňom dať pozor na to, ktorá os je vnútorná a~ktorá vonkajšia.

\EnvId{lHvoHwA3z301D2go6rtd5-existencia-stredu-vpisanej-kruznice}
\Theorem{existencia~$I$}{
    Vnútorné osi uhlov ľubovoľného trojuholníka sa pretínajú v~jednom bode. Tento je stredom kružnice vpísanej danému trojuholníku.
}{
    Nech $ABC$ je trojuholník. Nech~$I$ je priesečník osí vnútorných uhlov pri vrcholoch~$B$ a~$C$, ktorý zrejme existuje a~zrejme leží vnútri $ABC$.

    Podľa predošlej vety platí $d(I, AB) = d(I, BC)$ (lebo $I$ leží na osi uhla pri~$B$) a~$d(I, BC) = d(I, AC)$ (lebo $I$ leží na osi uhla pri~$C$). Odtiaľ $d(I, AB) = d(I, AC)$. Keďže~$I$ leží vnútri uhla~$BAC$, tak to znamená, že $I$ leží na jeho osi. Všetky tri vnútorné osi teda prechádzajú bodom~$I$.

    \Image{angles-incenter.pdf}

    \Remark{1} Bod~$I$ má rovnakú vzdialenosť od všetkých troch strán, takže je stredom kružnice, ktorá sa dotýka všetkých troch strán zvnútra -- teda kružnice vpísanej trojuholníku $ABC$.

    \Remark{2} Označenie~$I$ pochádza z~anglického slova \uv{incentrum}. V~česko-slovenskej olympiáde sa často používa proste~$S$ (ako stred).
}

V~úlohách sa veľmi často objavia aj \textit{kružnice pripísané}.

\EnvId{Mm527eX2RrorSGRJEbga4-existencia-stredu-pripisanej-kruznice}
\Theorem{existencia~$I_a$}{
    V~trojuholníku $ABC$ sa vnútorná os uhla pri~$A$ a~vonkajšie osi uhlov pri~$B$, $C$ pretínajú v~jednom bode. Tento je stredom kružnice pripísanej trojuholníku $ABC$ oproti vrcholu~$A$.
}{
    Pre prehľadnosť zápisu označme $X$, $Y$ body na polpriamkach opačných k~$CA$, $BA$. Vonkajšia os uhla pri~$B$ je práve vnútorná os uhla $YBC$ a~vonkajšia os pri~$C$ je vnútorná os uhla $XCB$. Označme~$I_a$ priesečník týchto osí.

    Z~vety o~osi uhla aplikovanej na uhly $YBC$ a~$XCB$ dostaneme rovnosti
    $$
    d(I_a, AB) = d(I_a, BC) = d(I_a, AC).
    $$
    Bod~$I_a$ má teda rovnakú vzdialenosť od priamok $AB$ a~$AC$, takže leží na vnútornej alebo vonkajšej osi uhla $BAC$. Zjavne však leží vnútri uhla $BAC$, takže ide o jeho vnútornú os -- tá teda prechádza bodom~$I_a$.

    \Image{angles-excenter.pdf}

    \Remark{1} Bod~$I_a$ má rovnakú vzdialenosť od priamky $BC$ aj~od priamok obsahujúcich strany $AB$, $AC$, takže je stredom kružnice dotýkajúcej sa strany $BC$ zvonku a~zvyšných dvoch strán (presnejšie ich predĺžení) zvnútra -- kružnice pripísanej trojuholníku oproti vrcholu~$A$.

    \Remark{2} Označenie~$I_a$ je bežné medzinárodné značenie. Kružnici pripísanej sa hovorí \uv{excircle}, takže možno logickejšie by mohlo byť $E_a$, ale to sa skôr nepoužíva. Kružnice vpísané ale súvisia s~incentrom, takže možno preto \UpsideDownSmile.
}

Osi uhlov, vnútorné i~vonkajšie, sú v~súťažných úlohách mimoriadne užitočný a~obľúbený koncept. Často sa skrývajú za iné podmienky a~kľúčový krok riešenia býva uvedomiť si, že nejaký bod leží na dvoch osiach, a~teda aj na tretej. Niekedy ide o~osi vnútorné, inokedy vonkajšie. Tie vonkajšie je často vidieť ťažšie. Určite sa s~osami oplatí kamarátiť.

\sec Čo si zapamätať

\secc Techniky

\begitems \style N
\i Pri rovnobežkách hľadáme zhodné striedavé a súhlasné uhly, prípadne uhly so súčtom~$180^\circ$.
\i Je užitočné nahliadať na vonkajší uhol trojuholníka ako na súčet dvoch uhlov.
\i Na prepojenie sveta dĺžok a~uhlov sa hodí zhodnosť dvoch trojuholníkov.
\i Pomocné body (stred úsečky, zrkadlový obraz, predĺženie) často odhalia zhodnosť alebo špeciálny trojuholník.
\enditems

\secc Užitočné fakty

\begitems \style N
\i Súčet uhlov v~trojuholníku je $180^\circ$, všeobecne v~konvexnom $n$-uholníku $(n-2) \cdot 180^\circ$.
\i Vety o~zhodnosti trojuholníkov: $sss$, $sus$, $usu$, $Ssu$.
\i V trojuholníku sú dve strany zhodné práve keď sú zhodné im protiľahlé uhly.
\i Pravouhlý trojuholník má preponu dvakrát dlhšiu než odvesnu práve keď má uhly $90^\circ, 60^\circ, 30^\circ$.
\enditems

\sec Úlohy

V~tejto sekcii nájdete rôzne úlohy, na ktoré vám stačia poznatky z~tohto materiálu, žiadne pokročilé veci ako \textit{obvodové a~stredové uhly} potrebné nie sú (na tieto úlohy sa pozrieme neskôr).


\EnvId{Vu5msR8YT8yuZ3zb8NeRJ-zhodne-usecky-v-trojuholniku}
\Problem{0}{MO okresné kolo Z8 2025}{
    V~trojuholníku $ABC$ leží bod $D$ na strane $BC$, bod $E$ na strane $AC$, pričom $|AB| = |BE| = |EC| = |CD|$ a~$|BD| = |DE|$. Určte veľkosti uhlov $ACB$ a~$BAD$.
}{
    Označte $\gamma = |\angle ACB|$ a~snažte sa všetky uhly v~obrázku vyjadriť cez~$\gamma$. Máme veľa rovnoramenností.
}{
    Jedna možná cesta je postupne pomocou $\gamma$ vyjadriť uhly $CBE$, $BED$, $CDE$, $CED$. Nevieme teraz niekde napísať rovnicu, z~ktorej sa dá vypočítať $\gamma$?
}{
    Pozrite sa na súčet uhlov v~trojuholníku~$BEC$. Z~neho by malo ísť spočítať $\gamma$.
}{
    K~dopočítaniu~$BAD$ potrebujeme viac krokov. Prvý je spočítať čo najviac uhlov a~niečo si všimnúť.
}{
    Keď spočítame uhly v~$ABC$, tak dostaneme, že $ABC$ je rovnoramenný. 
}{
    Teraz už vieme použiť symetriu, jedna možnosť je dokázať zhodnosť trojuholníkov $EAD$ a~$DBE$.
}{
    \Answer{$|\angle ACB| = 36^\circ$ a~$|\angle BAD| = 36^\circ$.}

    Označme $\gamma = |\angle ACB|$.

    V~rovnoramennom $\triangle BEC$ ($|BE|=|EC|$) je $|\angle EBC| = |\angle ECB| = \gamma$. Keďže $D$ leží na $BC$, je $|\angle DBE| = |\angle EBC| = \gamma$, a~z~rovnoramennosti $\triangle BDE$ ($|BD|=|DE|$) aj $|\angle DEB| = \gamma$. Veta o~vonkajšom uhle v~$\triangle BDE$ pri~$D$ dáva
    $$
    |\angle EDC| = |\angle DBE| + |\angle DEB| = 2\gamma,
    $$
    a~v~rovnoramennom $\triangle CDE$ ($|CD|=|CE|$) tak $|\angle CED| = |\angle CDE| = 2\gamma$. Zo súčtu uhlov v~ňom $\gamma + 4\gamma = 180^\circ$, čiže $\gamma = 36^\circ$.

    \Image{angles-isosceles-chain.pdf}{0.8}

    Teraz si vezmime, že $|\angle AEB|$ je vedľajší uhol k~$|\angle BEC|=3\gamma=108^\circ$, takže je rovný $72^\circ$. Z~rovnoramennosti $\triangle BAE$ tak aj $|\angle BAE|=72^\circ$. Z~$\triangle ABC$, v~ktorom už máme pri $C$ a~$A$ rovné postupne $36^\circ$ a~$72^\circ$ dopočítame $|\angle CBA|=72^\circ$, takže trojuholník $ABC$ je rovnoramenný.

    Teraz si vezmime, že z~$|CA|=|CB|$ a $|CE|=|CD|$ vlastne máme $|EA|=|DB|$. Vezmime si trojuholníky $EAD$ a $DBE$, tie sú rovnoramenné so zhodným uhlom pri vrchole (rovným $180^\circ-2\gamma$) a~zhodnými ramenami, takže podľa $sus$ sú zhodné. To dáva $|\angle DAE|=|\angle DBE|=36^\circ$. 

    Na záver teda máme
    $$
    |\angle BAD| = |\angle BAE| - |\angle DAE| = 2\gamma - \gamma = \gamma = 36^\circ.
    $$
}

\EnvId{SaOp1yOLr0T9xO-XywvHv-rovnostranne-trojuholniky-vo-stvorci}
\Problem{0}{DuoGeo 2025}{
    Do štvorca $ABCD$ boli nakreslené rovnostranné trojuholníky $ABX$ a~$CDY$. Určte súčet vyznačených uhlov.

    \Image{angles-square-equilateral-statement.pdf}{0.7}
}{
    Skúste najprv spočítať čo najviac uhlov na obrázku.
}{
    Kľúčom k~dopočítaniu finálneho uhla je nájsť vhodný rovnoramenný trojuholník cez rovnaké dĺžky.
}{
    \Answer{$180^\circ$.}

    Podľa symetrie sú všetky štyri vyznačené uhly zhodné. Stačí teda nájsť veľkosť jedného z~nich -- zameriame sa na $|\angle YDX|$.

    Najprv $|\angle XAD| = |\angle BAD| - |\angle BAX| = 90^\circ - 60^\circ = 30^\circ$; analogicky $|\angle ADY| = |\angle ADC| - |\angle YDC| = 90^\circ - 60^\circ = 30^\circ$.

    Ďalej $|DA| = |AB| = |AX|$ (prvé je strana štvorca, druhé strana rovnostranného trojuholníka $ABX$), takže $\triangle AXD$ je rovnoramenný so základňou $DX$. Z~rovnosti $|\angle XAD| = 30^\circ$ a~zo súčtu uhlov v~trojuholníku
    $$
    |\angle ADX| = \tfrac{1}{2}\bigl(180^\circ - 30^\circ\bigr) = 75^\circ.
    $$
    Hľadaný uhol je rozdielom dvoch už vypočítaných:
    $$
    |\angle YDX| = |\angle ADX| - |\angle ADY| = 75^\circ - 30^\circ = 45^\circ.
    $$
    Súčet všetkých štyroch vyznačených uhlov je teda $4 \cdot 45^\circ = 180^\circ$.

    \Image{angles-square-equilateral-solution.pdf}{0.8}
}

\EnvId{q5FKHnEWB5kfNu_ZGuq49-osi-uhlov-v-lichobezniku}
\Problem{0}{MO školské kolo C 2024}{
    V~lichobežníku $ABCD$, kde $AB \parallel CD$, sa osi vnútorných uhlov pri vrcholoch $C$ a~$D$ pretínajú na úsečke~$AB$. Dokážte, že $|AD| + |BC| = |AB|$.
}{
    Nech~$P$ je náš spoločný priesečník. Máme rovnobežnosť, máme os uhla, to dáva dosť rovnakých uhlov.
}{
    Cieľ je nájsť rovnoramenné trojuholníky.
}{
    Označme $P$ spoločný priesečník osí; podľa zadania leží na úsečke~$AB$. Ukážeme, že
    $$
    |AD| = |AP| \quad \hbox{a} \quad |BC| = |BP|,
    $$
    odkiaľ priamo $|AD| + |BC| = |AP| + |BP| = |AB|$.

    Pretože $DP$ je os vnútorného uhla pri~$D$, máme $|\angle ADP| = |\angle CDP|$. Z~rovnobežnosti $AB \parallel CD$ sú $|\angle CDP|$ a~$|\angle APD|$ striedavé uhly pri priečke~$DP$, takže
    $$
    |\angle APD| = |\angle CDP| = |\angle ADP|.
    $$
    V~trojuholníku $ADP$ sú teda uhly pri vrcholoch $D$ a~$P$ zhodné, takže $|AD| = |AP|$. Analogickou úvahou pri vrchole~$C$ dostaneme $|BC| = |BP|$.

    \Image{angles-trapezoid-bisectors.pdf}{0.8}
}

\EnvId{QyZrP8_I1-CPXFt7t7zed-rovnobeznost-v-konvexnom-patuholniku}
\Problem{0}{MO školské kolo A 2023}{
    V~konvexnom päťuholníku $ABCDE$ platí $|\angle CBA| = |\angle BAE| = |\angle AED|$. Na stranách $AB$ a~$AE$ existujú po rade body $P$ a~$Q$ tak, že $|AP| = |BC| = |QE|$ a~$|AQ| = |BP| = |DE|$. Dokážte, že $CD \parallel PQ$.
}{
    Máme veľa rovnakých uhlov a~strán, skúste nájsť zhodné trojuholníky.
}{
    Kľúčové sú zhodné trojuholníky $PBC$, $QAP$ a~$DEQ$. Z~nich dostaneme užitočné zhodné vecičky. Pokračujte v~hľadaní zhodnosti.
}{
    Finálna zhodnosť na dokázanie je $CPQ$ a $DQP$. Prečo to stačí? 
}{
    Keďže $|BC| = |AP| = |EQ|$, $|BP| = |AQ| = |ED|$ a~$|\angle CBP| = |\angle PAQ| = |\angle QED|$, sú podľa vety $sus$ trojuholníky $PBC$, $QAP$ a~$DEQ$ navzájom zhodné.

    \Image{angles-pentagon.pdf}{0.8}

    Odtiaľ plynie $|CP| = |PQ| = |QD|$ a~tiež
    $$
    \gather
    |\angle CPQ| = 180^\circ - |\angle BPC| - |\angle APQ| = \\ = 180^\circ - |\angle PQA| - |\angle EQD| = |\angle PQD|.
    \endgather
    $$
    Podľa vety $sus$ sú teda zhodné aj rovnoramenné trojuholníky $CPQ$ a~$DQP$. Z~toho vyplýva zhodnosť ich výšok z~vrcholov $C$ a~$D$ na spoločnú stranu $PQ$; tieto výšky sú zároveň rovnobežné (obe kolmé na $PQ$), takže $CD \parallel PQ$.    
}

\EnvId{my3-6U8Mr3piG2sLobH4A-dlzka-usecky-vo-stvoruholniku}
\Problem{0}{DuoGeo 2025}{
    Je daný štvoruholník $ABCD$ s~priesečníkom uhlopriečok $T$. Predpokladajme, že veľkosti uhlov $BAC$ a~$DBA$ sú postupne $30^\circ$ a~$45^\circ$. Na úsečke $BT$ leží bod $Z$ taký, že $CZ \perp BT$. Predpokladajme, že priamka $CZ$ pretne úsečku $AB$ v~bode $M$. Nech $R$ je priesečník úsečiek $AT$ a~$MD$. Predpokladajme, že $|AM| = |AR|$ a~$|MR| + |TD| = 14$. Určte veľkosť úsečky $|BZ|$.
}{
    Dopočítavajte uhly, až kým nenájdeme rovnoramenný trojuholník. 
}{
    Trpezlivým počítaním uhlov sa dá dôjsť k~tomu, že $DTR$ je rovnoramenný. To by malo vniesť svetlo do podmienky $|MR|+|TD|=14$.
}{
    Podmienka sa po dokázaní rovnoramenností preloží ako $|MD|=14$. To pomôže neskôr -- aktuálne už viac uhlov nespočítame a~musíme ešte niečo nájsť zo sveta dĺžok. Kľúčom je nájsť pekný pravouhlý trojuholník.
}{
    Dá sa spočítať, že uhly trojuholníka $MDZ$ sú $90-60-30$. Poznáme jeho preponu $MD$. Naše skvelé dokázané pomocné tvrdenie nám teraz dáva ďalšiu dĺžku. Potom je to krok od riešenia.
}{
    \Answer{$|BZ| = 7$.}

    V~trojuholníku $ATB$ poznáme uhly pri vrcholoch $A$ a~$B$, a~síce $30^\circ$ a~$45^\circ$. Tým pádom je vonkajší uhol pri~$T$ rovný súčtu týchto uhlov, konkrétne
    $$
    |\angle DTR|=|\angle TAB|=|\angle TBA|=30^\circ+45^\circ=75^\circ.
    $$
    V~rovnoramennom trojuholníku $AMR$ s~$|AM| = |AR|$ je uhol pri vrchole $A$ rovný $30^\circ$, takže oba uhly pri základni $MR$ majú veľkosť $\tfrac{1}{2}(180^\circ - 30^\circ) = 75^\circ$. Špeciálne $|\angle ARM| = 75^\circ$, a~teda aj jeho vrcholový uhol $|\angle DRT|$ má veľkosť $75^\circ$.

    \Image{angles-quad-30-45.pdf}{0.8}

    Trojuholník $DTR$ má teda dva uhly veľkosti $75^\circ$, čiže je rovnoramenný so~základňou $TR$ a~$|DT| = |DR|$. Odtiaľ
    $$
    |MD| = |MR| + |RD| = |MR| + |TD| = 14.
    $$
    Navyše $|\angle TDR| = 180^\circ - 2 \cdot 75^\circ = 30^\circ$.

    Pozrime sa teraz na trojuholník $MDZ$. Pretože $T$, $Z$ ležia na úsečke $BD$ a~$R$ leží na úsečke $MD$, je $|\angle MDZ| = |\angle RDT| = 30^\circ$. Trojuholník $MDZ$ je teda pravouhlý (pri~$Z$) s~preponou $MD$ a~uhlom $30^\circ$ pri~$D$, takže podľa skoršieho tvrdenia o~$30^\circ$-$60^\circ$-$90^\circ$ trojuholníku
    $$
    |MZ| = \tfrac{1}{2} |MD| = 7.
    $$

    Konečne trojuholník $MZB$ má pravý uhol pri~$Z$ a~uhol $45^\circ$ pri~$B$, takže aj $|\angle BMZ| = 45^\circ$. Je teda rovnoramenný pravouhlý a~$|BZ| = |MZ| = 7$.
}

\DisplayExerciseSolutions

\DisplayHints

\DisplayProblemSolutions

\bye